\section{Fundamentos teóricos}

\subsection{Filtro de Wiener}

O filtro de Wiener busca o filtro linear que minimiza o erro quadrático médio
entre a saída e um sinal desejado. Seja \(y = x + r\) a observação, em que
\(x\) é o sinal desejado e \(r\) o ruído. No domínio da frequência, para
processos estacionários de média nula, a resposta em frequência ótima é
\begin{equation}
H(e^{j\omega}) = \frac{S_{xy}(e^{j\omega})}{S_{yy}(e^{j\omega})},
\end{equation}
em que \(S_{yy}\) é a densidade espectral da observação e \(S_{xy}\) é o
espectro cruzado entre sinal desejado e observação. Quando o sinal desejado e o
ruído são incorrelacionados, \(S_{xy} = S_{xx}\) e
\(S_{yy} = S_{xx} + S_{rr}\), de modo que
\begin{equation}
H(e^{j\omega}) = \frac{S_{xx}(e^{j\omega})}{S_{xx}(e^{j\omega}) + S_{rr}(e^{j\omega})},
\end{equation}
em que \(S_{xx}\) é a densidade espectral do sinal desejado (limpo). No
presente trabalho, \(S_{xx}\) é modelada a partir do espectro do template,
escalado pela potência típica das amplitudes do treino, e \(S_{rr}\) é
estimada nas janelas de ruído puro na mesma convenção espectral, de modo que
as duas grandezas sejam diretamente comparáveis.

\subsection{Filtro casado}

O filtro casado maximiza a razão sinal-ruído no instante de interesse quando a
forma do sinal é conhecida. Para ruído branco, a resposta impulsiva é uma
versão revertida no tempo do template, \(h[n] = s[-n]\). A saída equivale à
correlação cruzada entre a observação e \(s[n]\). O máximo dessa saída é usado
como estatística de detecção, e a posição do máximo estima \(n_0\).

Quando o ruído é colorido, a formulação ótima envolve o branqueamento do ruído
(ou, de forma equivalente, o uso da covariância \(C^{-1}\)). Mesmo assim, o
filtro casado clássico permanece um detector simples e robusto e é o método
base adotado aqui para localização e detecção.

\subsection{Estimador BLUE da amplitude}

Fixado o instante \(n_0\), a observação na janela pode ser escrita como
\(y = A s_{n_0} + r\), em que \(s_{n_0}\) é o template alinhado em \(n_0\). O
melhor estimador linear não viesado (BLUE) da amplitude é
\begin{equation}
\hat{A} = \frac{s_{n_0}^{\mathsf{T}} C^{-1} y}{s_{n_0}^{\mathsf{T}} C^{-1} s_{n_0}}.
\end{equation}
Se o ruído for branco, \(C = \sigma^2 I\) e a expressão reduz-se a
\begin{equation}
\hat{A} = \frac{s_{n_0}^{\mathsf{T}} y}{s_{n_0}^{\mathsf{T}} s_{n_0}}.
\end{equation}
No trabalho, comparam-se as duas formas: BLUE branco e BLUE com a covariância
estimada a partir das janelas de ruído puro.

\subsection{Limiar de detecção e métricas}

Sob a hipótese \(H_0\) (apenas ruído), define-se um limiar \(\eta\) de modo que
a taxa de falso alarme seja aproximadamente \(1\%\). Em seguida, no conjunto
rotulado, avaliam-se:
\begin{itemize}
  \item reconstrução da forma: RMSE entre \(\hat{A}\,s[n-\hat{n}_0]\) e o pulso
  verdadeiro \(A\,s[n-n_0]\);
  \item detecção: curva ROC, AUC e, após o limiar, matriz de confusão, precisão,
  revocação e F1;
  \item amplitude: RMSE entre \(\hat{A}\) e \(A\) verdadeira nas janelas com
  evento.
\end{itemize}
